\section{Quantitative leakage from the actual exponential}\label{supp4:section:001}
Let $A$ be a unital complex C-star algebra, let $H=H^*$, and let $P=P^*=P^2$. Put $Q=1-P$. The checked statement is
\[
\boxed{\|Q\exp(-itH)P\|\le |t|\,\|HP-PH\|\qquad(t\in\R).}
\]
The matrix specialization uses the operator norm on the Hilbert space, implemented by \path{Matrix.Norms.L2Operator}. It is not an entrywise norm or the row-sum operator norm. The theorem includes the zero projector and empty matrix index type. The theorem also applies at time zero.

\subsection{Deriving the estimate by differentiation}
For a skew-adjoint generator $G$, define the actual conjugation path
\[
T_P(t)=\exp(tG)P\exp(-tG).
\]
The formalization derives its derivative from the exponential power series and the product rule:
\[
T_P'(t)=\exp(tG)(GP-PG)\exp(-tG).
\]
Both exponential factors are proved unitary. Left and right multiplication by a unitary preserve the C-star norm. Therefore
\[
\|T_P'(t)\|=\|GP-PG\|,
\qquad
\|T_P(t)-T_P(s)\|\le |t-s|\,\|GP-PG\|.
\]
The second inequality follows from the mean value inequality in normed spaces and covers either ordering of $s$ and $t$. It is not a finite-difference estimate or a numerical time discretization.

Multiplication by $\exp(tG)$ gives
\[
(T_P(t)-P)\exp(tG)=\exp(tG)P-P\exp(tG).
\]
For every orthogonal projector, $\|P\|\le1$ and $\|1-P\|\le1$. The exact corner identity
\[
(1-P)UP=(1-P)(UP-PU)P
\]
then bounds leakage by the exponential commutator. Taking $G=-iH$ gives the stated theorem because $\|[G,P]\|=\|[H,P]\|$.

\subsection{Checked consequences}
Unitarity also gives the capped estimate
\[
\|Q\exp(-itH)P\|\le\min\{1,\ |t|\,\|[H,P]\|\}.
\]
If $H=H_0+V$, both summands are self-adjoint, and $[H_0,P]=0$, then
\[
\|Q\exp(-it(H_0+V))P\|\le2|t|\,\|V\|.
\]
The latter uses $\|[V,P]\|\le2\|V\|$. Its constant is not claimed to be optimal. Exact commutation gives zero leakage. The work does not include unbounded generators or time-dependent Hamiltonians.



\section{Actual projected differentiation and curvature}\label{supp4:section:002}
Let $E,F$ be real normed spaces. Let $P:E\to\mathcal L(F,F)$ be differentiable with $P(x)^2=P(x)$ for every $x$. Let $s:E\to F$ be $C^2$ with $P(x)s(x)=s(x)$, and let $X,Y:E\to E$ be differentiable vector fields. The checked connection is the actual expression
\[
\nabla_Xs(x)=P(x)\,Ds(x)[X(x)].
\]
It is proved to preserve the range condition. It is linear in the vector field direction and obeys the scalar Leibniz rule
\[
\nabla_X(cs)=Dc[X]\,s+c\nabla_Xs
\]
for differentiable real scalar functions $c$. No differentiated projector identities are inserted as hypotheses in the main curvature theorem.

\subsection{Differentiated projector and section identities}
Differentiating $P^2=P$ and $Ps=s$ produces
\begin{align*}
DP[v]P+PDP[v]&=DP[v],\\
PDP[v]P&=0,\\
Ds[v]&=DP[v]s+P Ds[v].
\end{align*}
The zero sandwich follows algebraically from the first identity. All products preserve their noncommutative order.

\subsection{Retaining the bracket of vector fields}
Define the actual curvature action by
\[
R_P(X,Y)s=\nabla_X\nabla_Ys-\nabla_Y\nabla_Xs-\nabla_{[X,Y]}s.
\]
The source uses the Lie bracket defined through Fr\'echet derivatives. Symmetry of second derivatives is supplied by the $C^2$ hypothesis, not postulated separately. Expanding the connection gives
\[
R_P(X,Y)s=P\bigl(DP[X]Ds[Y]-DP[Y]Ds[X]\bigr).
\]
Substituting the section identity and removing the zero sandwich terms proves
\[
\boxed{R_P(X,Y)s=P\bigl(DP[X]DP[Y]-DP[Y]DP[X]\bigr)P s.}
\]
Thus the coefficient is precisely $P(\dd P\wedge\dd P)P$ with the convention
$(\dd P\wedge\dd P)(v,w)=DP[v]DP[w]-DP[w]DP[v]$.

\subsection{Tensor properties and the remaining global boundary}
The coefficient is proved bilinear, alternating, and supported on the range on both sides. The chain rule proves the pullback identity
\[
F_{P\circ f}(x)(v,w)=F_P(f(x))(Df(x)v,Df(x)w).
\]
The main theorem has global regularity hypotheses on normed parameter spaces. The pullback identity alone does not supply localization, smooth local trivializations, or gluing. The final supplement constructs the global range bundle and proves that the local connections glue.



\section{A constructive and total syndrome recovery}\label{supp4:section:003}
Let $P$ be an orthogonal code projector. The input certificate consists of a finite family $V_a$ satisfying
\[
V_aP=V_a,
\qquad
V_a^*V_b=\delta_{ab}P.
\]
These hypotheses state that the operators are partial isometries with orthogonal syndromes. The next supplement derives such a family from arbitrary Knill--Laflamme Gram data.

\subsection{Constructing the recovery on the whole ambient space}
Define the syndrome projectors and their complement by
\[
S_a=V_aV_a^*,\qquad S=\sum_aS_a,\qquad Q=1-S.
\]
The source proves self-adjointness, $S_aS_b=\delta_{ab}S_a$, $S^2=S$, and $Q^2=Q=Q^*$. It also proves $QV_a=0$ and $V_a^*Q=0$.

The recovery is not left undefined outside the syndrome subspaces. Its explicit Kraus operators are $R_a=V_a^*$ together with $R_{\mathrm{extra}}=Q$, giving
\[
\mathcal R(\rho)=\sum_aV_a^*\rho V_a+Q\rho Q.
\]
They satisfy exact completeness:
\[
\sum_aR_a^*R_a+R_{\mathrm{extra}}^*R_{\mathrm{extra}}
=\sum_aV_aV_a^*+Q=S+Q=I.
\]
For finite complex matrices, trace preservation and complete positivity at every finite auxiliary dimension are proved for this actual map. The amplified map is identified with the Kraus family $I_J\otimes R_a$, including the complement term. Positivity is not merely checked on a single matrix size.

\subsection{Exact correction, including coherent errors}
For every operator supported on the code, $P\rho P=\rho$, the central identity is
\[
\boxed{\mathcal R(V_a\rho V_b^*)=\delta_{ab}\rho.}
\]
Consequently normalized mixtures satisfy
\[
\mathcal R\!\left(\sum_ap_aV_a\rho V_a^*\right)=\rho
\quad\text{when}\quad\sum_ap_a=1.
\]
The identity is algebraic and is proved for complex weights. Physical mixtures use the nonnegative real case. For a coherent error $E_z=\sum_a z_aV_a$, the calculation of the cross terms yields
\[
\mathcal R(E_z\rho E_z^*)=\left(\sum_a z_a\overline{z_a}\right)\rho.
\]
Thus the same map also corrects normalized coherent superpositions. This construction supplies the recovery stage used in the full Knill--Laflamme proof that follows.



